\label{sec:dados}
A base econômica é a MIP brasileira de 2015, com 67 atividades e valores em milhões de reais a preços básicos \parencite{ibge2015mip}. A fronteira abrange emissões nacionais associadas à demanda doméstica e às exportações; exclui emissões estrangeiras. O destino identifica o setor do produto final, não seu comprador imediato ou país consumidor. A conta pela demanda não equivale à pegada mundial do consumo dos residentes.

As intensidades diretas de \textcite{sanguinetAzzoni2024}, em Gg de CO\textsubscript{2}/R\$ milhão de produção, são interpoladas entre 2011 e 2018, com pesos $3/7$ e $4/7$. Trata-se de uma aproximação: a evolução linear, a compatibilidade de preços e a alocação ambiental setorial limitam sua interpretação. A documentação de preparação e as verificações encontram-se no Apêndice~\ref{ap:metodo}; os resultados não representam um inventário físico validado de 2015.
